%%
%% OBJECTS/LOIS
%%
\section{Les lois:~la fonction de transition}

La \aconcept{fonction de transition} est l'\'el\'ement central de notre
probl\'ematique des environnements de programmation pour automates
cellulaires. %
Les objets pr\'esent\'es pr\'ec\'edemment ne sont pas particuli\`erement
sp\'ecifiques d'un environnement de d\'eveloppement d'AC\@. %
Ils pourraient convenir \`a un environnement de production et de
traitement d'images anim\'ees quelconques. %
Ils constituent des extensions plus ou moins \'evidentes d'outils
existants\footnote{%
  Initialement nous pensions m\^eme pouvoir utiliser un des outils
  existants~(dans le domaine de l'animation interactive d'images),
  auquel des fonctions sp\'ecifiques d'application de fonctions de
  transitions auraient pu \^etre ajout\'ees.} %
ou sont bas\'es sur des imp\'eratifs d'architectures et de gestion de
ressources totalement in\'evitables. %

La fonction de transition concentre les probl\`emes de gestion de
ressources en m\'emoire et en CPU, et le probl\`eme du pouvoir expressif
de la sp\'ecification du r\'eseau d'interconnexion cellulaire et de la
table de transition. %

\mydumpfullR[t]{.5}{loop-1024}%
{Un espace produit par une table import\'e}%
{Une table import\'ee pour la g\'en\'eration~1024 de l'automate auto-reproducteur de Langton}%
{la table de la r\`egle~\arule{loop} qui a produit cette image vient de
  \asoft{cellsim}. %
  Le source servant \`a produire la table ne faisant pas partie de
  \asoft{cellsim}\footnote{%
    La reconstruction d'un tel source ne pose bien s\^ur pas de
    probl\`eme, mais nous souhaitions user des propri\'et\'es des tables.} %
  Un examen sommaire du code source de \asoft{cellsim} permet de
  d\'eterminer l'ordre choisi pour l'index, et un programme de
  permutation des entr\'ees de la table en fonction de cet ordre permet
  la conversion de la table au format de notre moteur.}

\myfigtwothirdsR[t]{.75}{func-1}%
{}%
{Deux visions du voisinage}%
{Cette figure montre la transformation d'un fragment
  d'espace en index d'une table de transition. %
  Le choix du mapping des coordonn\'ees des cellules de l'espace
  (de dimension $d$) vers le vecteur d'index est arbitraire. %
  Ce choix n'est pas neutre cependant, car il peut \^etre d\'eterminant pour
  l'efficacit\'e de la fonction de mapping. %
  \mygetfig[phrase]{apply-1}}

\myfigfull[t]{apply-1}%
{Le pipeline d'application}%
{Trois \'etapes successives de collecte de voisinage, illustrant le
  principe du pipeline de collecte mis en {\oe}uvre par la fonction
  d'application.}%
{Cette figure montre comment pipeliner la fonction de mapping pour un
  voisinage de Moore, r\'eduisant l'obtention d'un nouvel index \`a
  1~shift et 3~copies, \`a la place des 9~copies
  n\'ecessaires en l'absence de pipeline. %
  \mysetref{mem-struct}}

\myfigtwothirds[t]{apply-2}%
{}%
{Quelques r\'eseaux d'interconnexion}%
{Quelques r\'eseaux d'interconnexions. %
  Les premier et troisi\`eme, ainsi que les deuxi\`eme et troisi\`eme
  d'entre eux, dont les bits cibles sont disjoints, sont applicables
  en parall\`ele sur le m\^eme espace. %
  Les deuxi\`eme et troisi\`eme, de tailles d'index et de tailles de valeur
  cible identiques, peuvent \^etre associ\'es \`a une table de transition
  identique.}

\myfigtwothirds[t]{func-2}%
{}%
{Mise \`a jour par table}%
{Cette figure montre l'application d'un r\'eseau d'interconnexion sur
  une table.}

\myfigtwothirds[t]{func-3}%
{}%
{Deux tables disjointes}%
{Mise en {\oe}uvre de deux AC disjoints partageant un m\^eme espace
  d'exploitation.}

\myfigtwothirds[t]{func-4}%
{}%
{Deux tables non disjointes}%
{Mise en {\oe}uvre de deux AC coordonn\'es par l'interm\'ediaire de deux
  r\'eseaux d'interconnexions non disjoints partageant un m\^eme espace
  d'exploitation, permettant la transmission d'information d'un plan
  cellulaire \`a l'autre.}

\myfigfull[t]{func-7}%
{Fonction d'application pour un voisinage sp\'ecial}%
{Une fonction d'application pour un voisinage sp\'ecial dot\'e d'un
  accumulateur additif.}%
{Une fonction d'application pour un voisinage de Moore o\`u l'addition
  des cellules sources remplace leur concat\'enation. %
  Les connexions barr\'ees sont activ\'ees au cycle~2, et les connexions
  doublement barr\'ees au cycle~3 du pipeline.}

%%
%% OBJECTS/LOIS/STRUCTURE
%%
\subsection{Structure des lois:~voisinage et lookup-table}

La fonction de transition est la structure de contr\^ole de l'AC\@. %
Une fonction de transition locale re\c{c}oit en argument un fragment
d'espace~(un ensemble de cellules ordonn\'ees)\footnote{%
  Encore appel\'ee voisinage.}. %
Il existe un tel fragment pour chaque cellule de l'espace. %
La fonction retourne le nouvel \'etat de la cellule associ\'ee \`a ce
fragment~(la cellule cible). %
En g\'en\'eral, la cellule cible est un des \'el\'ements du fragment. %
Si le fragment n'est pas r\'eduit \`a la cellule cible, les fragments ne
sont pas disjoints et l'information peut circuler de cellule en
cellule. %

Sur une machine s\'equentielle, un espace temporaire stocke le nouvel
\'etat des cellules cibles et l'espace temporaire devient le nouvel \'etat
d'espace apr\`es l'application de la fonction de transition locale \`a
chaque cellule de l'espace courant\footnote{%
  Avec la contribution du s\'equenceur.}. %
Ce m\'ecanisme de simulation d'une architecture SIMD dot\'ee d'un r\'eseau
d'interconnexion homog\`ene ne pr\'esente pas de difficult\'e algorithmique
mais l'efficacit\'e et la souplesse du syst\`eme d\'ependent de
l'architecture choisie pour sa r\'ealisation. %

Il faut une m\'ethode de \aconcept{balayage} et de \aconcept{collecte}
des fragments d'espace pour chaque cellule cible,
d'\aconcept{association}\footnote{%
  Encore appel\'e liaison ou \aquote{binding}.} %
des param\`etres formels de la fonction de transition aux cellules du
fragment et d'\aconcept{invocation} de la fonction. %
La fonction qui re\c{c}oit en argument un ensemble d'espaces, un r\'eseau
d'interconnexion et une fonction de transition, est la
\aconcept{fonction d'application}\footnote{%
  Encore appel\'ee moteur cellulaire.}. %

Le choix du m\'ecanisme d'invocation conditionne la m\'ethode de liaison. %
La \emph{m\'ethode} la plus \emph{rapide} d'\emph{invocation} de
\emph{fonction} est l'acc\`es par \emph{table}. %
On construit une table des valeurs de la fonction pour l'ensemble des
valeurs possibles de ses arguments et chaque \'el\'ement de l'ensemble des
valeurs est consid\'er\'e comme un indice d'acc\`es \`a la table. %
La liaison des valeurs d'appel de la fonction se ram\`ene \`a la
concat\'enation des valeurs des cellules du fragment d'espace et le
calcul du nouvel \'etat de la cellule cible \`a un acc\`es m\'emoire
indirect. %

Les avantages sont \'evidents. %
Les fonctions de transitions peuvent \^etre r\'ealis\'ees dans n'importe
quel langage, sans souci de performance. %
L'op\'eration de mise \`a jour d'une cellule est r\'eduite \`a un nombre
minimum d'instructions. %
Les inconv\'enients sont \'egalement \'evidents. %
Les fonctions r\'ealisables sont limit\'ees par la taille des tables~(en
nombre d'entr\'ees) qui est une fonction exponentielle de la taille des
arguments~(produit du nombre de bits~$k = \log_{2}(K)$ codant le
nombre~$K$ d'\'etats par cellules et du nombre~$r$ de cellules du
fragment d'espace). %
La taille~$T$ des tables exprim\'ee en bytes d\'epend de~$k$ exprim\'e
en~bytes, $T = \byte(k) 2^{kr}$. %
Le rendement de la contraction de la table au strict nombre de bits
n\'ecessaires est proportionnel au rapport du nombre de voisins et du
nombre d'\'etats. %
Un espace de dimension~$d = 3$ dot\'e du voisinage complet %
des $r = 26$~cellules distantes de 1~unit\'e de la cellule cible
\mygetref[paren]{nombre-de-voisin} et de $k = 1$~bit d'\'etat requiert
une table de 128~m\'egabytes sans contraction et 16~m\'egabytes avec
contraction. %

Nous pensons que les avantages sont plus grands que les inconv\'enients. %
Le nombre d'AC accessibles, m\^eme limit\'e, reste \'enorme, et l'\'etude des
AC concerne pour une bonne part les plus simples d'entre eux. %
De plus il est possible de cr\'eer des fonctions complexes
en combinant\footnote{%
  \mygetref[phrase]{combine-func}} %
l'application de tables de fonction simple. %
%% Cette id\'ee constitue en fait le c{\oe}ur de notre recherche,
%% et la seule v\'eritable source de cr\'eation d'objets cellulaires.

%%
%% OBJECTS/LOIS/OPERATIONS
%%
\subsection{Op\'erations:~plasticit\'e des tables de transitions}

Un autre avantage des tables est de permettre tr\`es simplement la
g\'en\'eration ou la transformation de fonctions. %
Un param\`etre \'el\'ementaire de g\'en\'eration de fonctions est par exemple
donn\'e par la densit\'e de remplissage de la table\footnote{%
  Il s'agit du param\`etre $\lambda$~\cite{Langton92b}, d'une richesse
  \'etonnante~\mygetref[paren]{lambda} puisqu'il est directement associ\'e
  \`a la classe de comportement dynamique de la fonction.}. %
En fait, une fonction assimil\'ee \`a un vecteur d'\'etats devient un objet
cible de multiples r\'eductions et transformations au m\^eme titre qu'un
espace de dimension~1. %
Des attributs comme la tranquillit\'e\footnote{%
  Un voisinage dont toutes les cellules sont \`a z\'ero~(un \'etat
  conventionnel nul), donne une cellule cible \'egalement \`a z\'ero.} %
ou l'isotropie impliquent, cependant, une connaissance du mapping
\bialt{espace}{index} du voisinage. %
Le simple comptage des acc\`es aux entr\'ees de la table pour une
transition~(ou cumul\'e sur le temps) fournit une information
statistique relativement complexe\footnote{%
  L'histogramme de l'occurrence de toutes les combinaisons d'espaces
  possibles sur un fragment d'espace de la taille du voisinage,
  qui peut servir \`a calculer l'entropie de l'espace.} %
sur le point d'espace pr\'ec\'edent pour un co\^ut CPU
n\'egligeable\footnote{% 
  Au prix n\'eanmoins d'une consommation m\'emoire~$M$ p\'enalisante pour
  les plus grandes tables %
  $M = \log_{2}(SN)T$, en pratique $2 T$ ou $4 T$ bytes.}. %
On peut encore citer les op\'erateurs de bruitage~(ou mutation) sur
table et de cross-over sur deux tables, qui permettent la r\'ealisation
d'algorithmes g\'en\'etiques pour l'exploration de l'espace des fonctions,
ou enfin l'utilisation possible de r\'eseaux de neurones sur une
table. %

%%
%% OBJECTS/LOIS/COMPOSITION
%%
\subsection{Composition:~les chemins de donn\'ees}

\mysetref{combine-func}%
La coordination de plusieurs tables n\'ecessite la mise au point d'un
syst\`eme de liaison \bialt{param\`etres}{valeurs} pour le multiplexage spatial
des tables et d'un s\'equenceur pour le multiplexage temporel. %

La n\'ecessit\'e d'adresser rapidement chaque cellule de l'espace sur des
machines adressables au niveau du byte, conduit au choix d'une
organisation en plan de vecteur de bits~(tableau de bytes) plut\^ot
qu'en vecteur de plan de bits. %
Le r\'eseau d'interconnexion est une liste de couples
\bialt{cellule}{bit} sp\'ecifiant les coordonn\'ees spatiales de chacun
des bits constituant l'index d'acc\`es \`a la table, associ\'e \`a une liste
des bits cibles de la valeur de l'entr\'ee dans la table. %

\myfigfull[t]{gap-pe}%
{}%
{Structure d'un PE GAPP}%
{Reproduite d'apr\`es~\cite{NCR85a}, sans faire figurer le registre CM}%

\myfigfull[t]{gap}%
{}%
{Sch\'ema de la fonction d'application de la table \arule{gapp}}
{Les fl\`eches figurant sur le sch\'ema indique les chemins de donn\'ees des
  bits constituant l'index de la table \arule{gapp}. %
  Les conditions d'affectation et les temps de cycles internes \`a la
  fonction d'application sont indiqu\'es \`a droite du sch\'ema.}

\mycodefull[t]{c-gapp-pe-h}%
{}%
{La structure de donn\'ees de description de l'architecture d'un PE GAPP}%
{La simplicit\'e de l'impl\'ementation appara\^it clairement. %
  Les champs \avar{Bit\_Field} contiennent les parties d\'eclaratives de
  l'architecture. %
  Les deux m\'ethodes associ\'ees \`a cet objet sont la fabrication de la
  table de transition et la fabrication de la table de s\'election d'une
  fonction d'application qui r\'ealise l'association
  \bialt{instruction}{fonction d'application}.}

\mycodefull[t]{c-gapp-pe-pe}%
{}%
{La fonction impl\'ementant un PE GAPP}%
{Lib\'er\'e des contraintes h\'erit\'ees du m\'elange de la sp\'ecification et de
  l'ex\'ecution, la fonction impl\'ementant le PE GAPP est un mapping
  presque imm\'ediat de sa sp\'ecification architecturale\footnote{%
    L'ajout d'une sp\'ecification de m\^eme type pour la partie ALU,
    permettra la synth\`ese automatique de fonctions d'impl\'ementation de
    PE \`a partir de sp\'ecification purement d\'eclaratives.}.}

\mycodefull[t]{c-gapp-pe-tbl}%
{}%
{La fabrication de la table de transition de l'AC \aquote{PE GAPP}}%
{La fabrication de la table ne pose \'evidemment aucun probl\`eme, il
  s'agit d'une simple boucle. %
  La figure fait appara\^itre clairement comment le passage par table
  lib\`ere le style de l'impl\'ementation, en effet les fonctions
  \avar{implode} et \avar{explode} qui assurent les conversions entre une
  repr\'esentation bit-packed et bit-expanded\footnote{%
    Dans laquelle un bit est stock\'e dans un mot de la taille naturelle
    de l'architecture, ici un entier de 32~bits est utilis\'e pour
    stocker 1~bit.} %
  peuvent \^etre utilis\'ees sans pr\'eoccupations inu\-tiles de
  performances ou d'occupation m\'emoire\footnote{%
    La pr\'esence d'une d\'eclaration \aquote{inline}, indique toutefois
    qu'il reste possible d'utiliser des optimisations qui ne
    compromettent pas la g\'en\'eralit\'e de l'impl\'ementation. %
    En phase de d\'eveloppement, la rapidit\'e de construction de la table
    est importante.}.}

\mycodefull[t]{c-gapp-arch-h}%
{}%
{La structure de donn\'ees de la partie d\'eclarative de la description de
  l'architecture d'un PE GAPP}%
{Un PE est compos\'e d'un champ de bits pour les entr\'ees, les sorties,
  l'instruction, et d'une table de v\'erit\'e pour l'alu\footnote{%
    Cette table est ici donn\'ee in extenso sans description
    associ\'ee~\mygetcode[paren]{c-gapp-arch-alu}.}. %
  Une description purement d\'eclarative d'une architecture de PE,
  n\'ecessiterait la mise en champ de bits de celle-ci.}

\mycodefull[p]{c-gapp-arch}%
{}%
{Les initialisation de la structure de donn\'ees de la partie d\'eclarative
  de la description de l'architecture d'un PE GAPP}%
{On conviendra ais\'ement de la possibilit\'e d'une g\'en\'eration
  automatique d'un telle description \`a partir d'un syntaxe de plus
  haut niveau.}

\mycodefull[t]{c-gapp-arch-alu}%
{}%
{Initialisation de l'ALU d'un PE GAPP}%
{Une architecture mixte \bialt{SIMD}{cellulaire} consid\`ere l'ALU comme
  une table de transition d'un AC, et permet le s\'equencement des
  chemins de donn\'ees \`a travers ces tables.}

%%
%% OBJECTS/LOIS/VOISINAGES-SPECIAUX
%%
\subsection{Voisinages sp\'eciaux:~n\'ecessaire compromis ou g\'en\'era\-lisation?}

On remarque sur la figure~\mygetfig{func-4} la pr\'esence d'un
\'el\'ement~(X) de voisinage non issu de l'espace. %
La fonction d'application peut fournir une valeur fonction du temps ou
de l'espace utilisable comme \'el\'ement de l'index. %
Ce type d'extension \`a la notion de voisinage peut sembler contraire \`a
la d\'efinition que nous avons donn\'e d'un AC\@. %
M\^eme si les AC qui nous int\'eressent sont b\^atis par combinatoire d'AC
simples, il existe toujours un couple \aquote{espace initial, fonc\-tion
  de transition} permettant de calculer et de diffuser une valeur
globale quelconque de l'espace ou du temps\footnote{%
  Ou si l'on pr\'ef\`ere il est possible de simuler une architecture de
  von Neuman dans un espace cellulaire.}. %
Les \'el\'ements de voisinages sp\'eciaux sont des acc\'el\'erateurs d'acc\`es \`a
ces valeurs. %
Certains d'entre eux sont simplement des \'economiseurs d'espace. %
Par exemple, au lieu d'utiliser un plan m\'emoire pour stocker un
pattern r\'egulier permettant \`a l'AC de conna\^itre la parit\'e des
coordonn\'ees d'une cellule\footnote{%
  Ce qui revient \`a organiser un multiplexage$^*$ spatial des r\`egles,
  qui pourrait \'egalement \^etre r\'ealis\'e en acc\'edant \`a la table de
  transition \`a travers une indirection suppl\'ementaire via un index de
  type spatial d'une cellule. $^*$Param\'etrage des r\`egles selon leurs
  positions dans l'espace.}, %
la fonction d'application maintient un bit de parit\'e par dimension
d'espace. %
D'autres permettent de r\'eduire la taille d'une table en fournissant
un compteur ou un timer synchrone sur l'espace\footnote{%
  Dans ce cas, la table contient plusieurs fonctions de transition
  distinctes qui alternent cycliquement, et il est sans doute plus
  avantageux d'utiliser un s\'equenceur externe.}. %
Enfin certaines r\'eductions de l'espace, dont le calcul n\'ecessiterait
un protocole cellulaire complexe, sont \'evidemment utiles, tel le OR de
l'ensemble des valeurs d'un plan de bit qui permet \`a chaque cellule de
savoir que ce plan est \`a z\'ero pour toutes les cellules de l'espace. %

La fonction d'application peut \'egalement calculer une valeur locale,
permettant ainsi une r\'ealisation par tables de fonctions autrement
inaccessibles. %
Si la fonction de transition op\`ere une r\'eduction simple sur l'index,
celle-ci peut \^etre confi\'ee \`a la fonction d'application. %
Si l'op\'erateur de r\'eduction est commutatif, associatif, et poss\`ede un
op\'erateur inverse, la r\'eduction peut \^etre pipelin\'ee. %
Par exemple, un voisinage de $r$ cellules \`a $k$ bits d'\'etat
donnerait une table de taille $2^{rk}$. %
La distribution de l'op\'erateur d'addition sur les $r$ cellules
sources permet de r\'eduire la taille de la table \`a %
%%$2^{\log_{2}(r 2^{k})}$. %
$r 2^{k}$. %
Dans ce cas pour $k = 8$, la taille de la table passe de %
$2^{72}$~bytes \`a $2^{12}$~bytes, et le pipeline de liaison remplace une
op\'eration de d\'ecalage\footnote{%
  A condition de disposer d'une arithm\'etique sur mots de 72 bits.} %
par une soustraction, et trois concat\'enations~(un shift et un or) par
trois additions. %

La distribution de l'op\'erateur d'addition sur les $r - 1$~cellules
sources (priv\'ees de la cellule cible dont la pr\'eservation est en
g\'en\'eral souhaitable pour obtenir une gamme d'AC accessibles plus
vaste) associ\'ee \`a la concat\'enation de la cellule cible,
donnerait une taille
%%2^(log((R - 1) * 2^K) + K)
%%2^(log(R - 1) + K + K)
%%2^(log(R - 1) + 2K)
$M = 2^{\log_{2}((r - 1) 2^{k}) + k} = 2^{\log_{2}(r - 1) + 2k} = (r-1)2^{k}$. %
Cette architecture op\`ere l'addition pipelin\'ee des cellules du
voisinage de Moore moins la cellule cible, et la concat\'enation de
cette somme avec la valeur de la cellule cible. %
Dans ce cas, pour $k = 8$, la taille de la table passe \`a~$2^{19}$~bytes. %

D'autres r\`egle locales, telles les minimums et maximums des
valeurs cellulaires, pourraient \^etre calcul\'ees par la fonction
d'application. %
L'association d'AC dont les cellules poss\`edent une vari\'et\'e de densit\'e
large pour une vari\'et\'e de comportements r\'eduite, \`a des AC de nature
oppos\'ee, est un de nos objectifs. %

En plus d'\'el\'ements de voisinages sp\'eciaux, la fonction d'application
peut fournir toutes sortes de voisinages non standards, destin\'es \`a une
utilisation dans le cadre du s\'equenceur, ou \`a un usage g\'en\'eralis\'e du
m\'ecanisme des tables de transitions. %
Nous avons d\'ej\`a \'evoqu\'e~\mygetref[paren]{time-xor} un voisinage \`a deux
bits associ\'e \`a la table \arule{Xor} pour la r\'ealisation d'une transformation
sur le temps. %
Le s\'equencement d'une s\'erie de voisinages \`a deux cellules, balayant un
espace r\'egulier standard en pivot sur la cellule cible, permet
d'accumuler une op\'eration quelconque sans faire appel \`a un \'el\'ement de
voisinage sp\'ecial~(non spatial) pour cette
r\'eduction\footnote{%
  \mysetref{seq-vs-pipe}%
  Le nombre d'op\'erations n\'ecessaires \`a la fabrication d'un index est
  r\'eduit mais le balayage de la m\'emoire est multipli\'e et la pr\'esence
  possible d'un pipeline pour les voisinages disposant d'un plus grand
  nombre de cellules r\'eduit cet avantage.}. %
Une table \`a deux arguments, utilisable comme accumulateur, a deux
entr\'ees de tailles diff\'erentes et une sortie de taille \'egale \`a
l'entr\'ee accumulatrice. %
Cette table peut \^etre utilis\'ee par une fonction d'application pour le
calcul d'une valeur locale en place de la concat\'enation pour la
fabrication de l'index. %

Dans le cas d'une g\'en\'eralisation de l'usage de tables \`a une
combinatoire~(multiplexage spatial et temporel) d'usage de tables non
disjointes sur un m\^eme espace, le gain escompt\'e n'est pas toujours un
gain de m\'emoire, mais \'egalement du pouvoir d'expression des structures
de donn\'ees associ\'ees \`a cette combinatoire. %
Par exemple, il n'est pas n\'ecessaire de consid\'erer le OR de tous les
bits d'un plan cellulaire comme un \'el\'ement de voisinage sp\'ecial dont
le calcul est laiss\'e au soin de la fonction d'application. %
Il est possible d'utiliser une variable sp\'eciale, un accumulateur~(une
variable scalaire, contrairement aux variables cellulaires), comme un
\'el\'ement de cellule cible, et non plus seulement d'index. %
Le binding de la table de la fonction binaire OR, ou de toute autre
fonction binaire, aux arguments CENTRE et ACU et au r\'esultat ACU, et
le report au s\'equenceur de la t\^ache de l'application de cette fonction
sur l'espace~(ou plus probablement \`a la capacit\'e de la fonction
d'application d'appliquer plusieurs AC sur un m\^eme pipeline) permet
d'utiliser la machinerie cellulaire \`a son propre service. %
Bien s\^ur, dans ce cas, seul le r\'esultat de la r\'eduction appliqu\'ee sur
le plan pr\'ec\'edent est accessible~(\`a moins de faire deux passes), et ce
type d'approche n'est pas possible sur une architecture SIMD r\'eelle. %

%%
%% OBJECTS/LOIS/UTILISATION
%%
\subsection{Utilisation:~construction, composition, m\'elange et
  s\'e\-paration}

Si l'on veut consid\'erer l'application de tables sur des espaces,
lesquels ne sont pas toujours confondus avec l'espace cellulaire \`a
proprement parler, comme la structure de contr\^ole
\'el\'ementaire\footnote{%
  En tout cas dans un environnement SISD\@. %
  Nous choisissons d\'elib\'er\'ement d'ignorer les probl\`emes que poserait
  le portage du syst\`eme sur une architecture SIMD\@. %
  Les AC \'etant quant \`a eux \'evidemment ind\'ependants de
  l'architecture. %
  Nous pensons que l'usage d'une architecture SIMD simul\'ee reposant
  sur une organisation pipelin\'ee de la transformation de plans
  m\'emoires par tables permet une utilisation efficace des ressources
  d'une architecture von Neumann, applicable \`a une vaste classe de
  probl\`emes g\'en\'eraux.} %
de l'univers cellulaire, celle-ci doit \^etre munie d'une structure de
donn\'ees supportant la construction, le m\'elange et la composition. %

%%
%% OBJECTS/LOIS/USAGE/CONSTRUCTION
%%
\begin{itemize}
\item%
  \mysetref{static-param}%
  La construction des tables permet l'association d'une fonction \`a ses
  param\`etres statiques. %
  Les param\`etres statiques d'une fonction sont les param\`etres de la
  fonction qui ne sont pas associ\'es \`a un des \'el\'ements du voisinage
  cellulaire, mais permettent de g\'en\'eraliser le comportement de la
  fonction. %
  Il s'agit par exemple de la valeur de timers locaux~(propres \`a
  chaque cellule) ou de valeurs de seuil d\'eclenchant des portions de
  code de la fonction. %

%%
%% OBJECTS/LOIS/USAGE/COMPOSITION
%%
\item%
  La composition des tables permet la construction d'une liste de
  fonctions travaillant en parall\`ele sur un m\^eme espace. %
  On peut obtenir le m\^eme effet en utilisant un s\'equenceur temporel,
  mais au d\'etriment de l'usage du pipeline\footnote{%
    \mygetref[phrase]{seq-vs-pipe}}. %
  De plus, la disponibilit\'e d'un objet sp\'ecifique de composition de
  fonctions permet l'application d'op\'erateurs de m\'elange et de
  s\'eparation sur les tables. %
  On peut enfin imaginer une g\'en\'eralisation de cette liste \`a un arbre,
  permettant l'utilisation de r\'esultats interm\'ediaires d'application
  de table comme \'el\'ement de voisinage d'une autre table. %

%%
%% OBJECTS/LOIS/MELANGE
%%
\item%
  Le m\'elange de deux tables consiste en la combinaison de leurs
  param\`etres formels en une seule liste de param\`etres, r\'esultant en
  une table dont la taille est \'egale au plus~(si l'intersection des
  ensembles de param\`etres est vide) au produit de la taille de chaque
  table\footnote{%
    Il s'agit de la taille de la table en nombre d'entr\'ees, qui doit
    \^etre multipli\'ee par le cardinal de l'union des param\`etres,
    exprim\'e en bits, pour obtenir la taille en bits.}. %
  Cette op\'eration est applicable soit au niveau des fonctions elles
  m\^emes, et n\'ecessite une nouvelle invocation de chaque fonction
  pour la construction de la nouvelle table\footnote{%
    Cela peut s'av\'erer relativement co\^uteux pour les grandes tables,
    si la fonction est fournie par un interpr\`ete.}, %
  soit directement sur les tables elles m\^emes. %
  Le m\'elange de tables permet une optimisation de l'usage des
  ressources CPU au d\'etriment des ressources m\'emoire. %
  Cette optimisation peut sembler douteuse puisqu'elle \'echange une
  augmentation lin\'eaire\footnote{%
    Le facteur de gain d\'epend de fa\c{c}on difficile \`a pr\'evoir du co\^ut
    relatif de la construction s\'epar\'ee ou unifi\'ee de l'index et de la
    dispersion des index produit par la simulation.} %
  de la vitesse de calcul contre une augmentation exponentielle de la
  taille de la m\'emoire utilis\'ee. %
  De plus la construction de grandes tables redondantes est en
  elle-m\^eme un processus co\^uteux, voire tr\`es co\^uteux si l'on en vient
  \`a faire appel \`a la m\'emoire pagin\'ee. %
  Cependant une m\'emoire qui se compte en dizaine de m\'egabytes sur une
  station de travail, rend int\'eressant ce genre d'optimisation. %
  D'autant plus qu'une fois achev\'ee la construction de la table,
  son usage peut r\'ev\'eler une densit\'e d'acc\`es concentr\'ee sur des
  fractions de la table\footnote{%
    Cela d\'epend des espaces initiaux et des AC, bien s\^ur. %
    Tous ceux qui ont observ\'e \arule{life} partant d'un espace
    initial al\'eatoire de densit\'e\,\protect\Fundemi,
    comprennent qu'apr\`es les premiers points de temps
    la densit\'e d'impact des index se concentre en fonction de la
    densit\'e g\'en\'erale de l'espace,
    et de la densit\'e des patterns produits.} %
  allant jusqu'\`a justifier l'usage de tables d'une taille
  sup\'erieure \`a la m\'emoire r\'eelle disponible. %
  Encore que dans ce cas, le cache devrait plut\^ot \^etre fourni par
  l'applicatif, sous forme d'une construction virtuelle de la
  table\footnote{%
    Pour ne pas \^etre trop co\^uteuse, celle-ci doit reposer sur les
    services du syst\`eme d'exploitation qui permettent une
    r\'ecup\'eration~(trap) des d\'efauts de pages.}, %
  que par le syst\`eme d'exploitation. %
  La dualit\'e opposant l'usage de nombreuses petites tables appliqu\'ees
  en parall\`ele sur un m\^eme espace, \`a l'usage d'une seule grande ou
  tr\`es grande table, permet de choisir entre une mise en~{\oe}uvre
  adapt\'ee au d\'eveloppement incr\'emental interactif~(faible d\'elai de
  mise en marche d'une architecture, associ\'e \`a un d\'ebit de mise \`a jour
  cellulaire non optimum), et une mise en~{\oe}uvre adapt\'ee \`a
  l'exploration de la dynamique \`a long terme d'AC sur de grands
  espaces~(temps de d\'emarrage important, associ\'e \`a un d\'ebit de mise \`a
  jour cellulaire optimum pour une architecture mat\'erielle donn\'ee). %

%%
%% OBJECTS/LOIS/SEPARATION
%%
\item%
  La s\'eparation des tables est l'op\'eration inverse du m\'elange. %
  Cette op\'eration~(pour laquelle nous n'avons pas actuellement cherch\'e
  d'algorithme) devrait permettre de d\'eterminer la qualit\'e de
  primitive d'une fonction. %
\end{itemize}

\myfigfull[t]{life-dfa}%
{}%
{Le DFA de \arule{Life}}%
{La construction automatique du DFA de nombre d'\'etats minimum associ\'e
  \`a une table ne pose pas de probl\`eme:~il suffit de construire l'arbre
  de d\'ecodage des index, et d'appliquer sur cet arbre consid\'er\'e
  comme un DFA un algorithme de minimisation.}

\myfigfull[t]{life-tree}%
{}%
{Une partie de l'arbre de d\'ecodage de \textbf{life}}%
{les 32 premi\`eres entr\'ees de la table de transition. %
  La racine de l'arbre est l'\'etat initial du DFA, les feuilles de
  l'arbre sont les entr\'ees de la table. %
  Une entr\'ee \`a~1 est un \'etat acceptant du DFA\@.}

%%
%% OBJECTS/LOIS/VIRTUAL
%%
\subsection{Des machines virtuelles:~compilation d'architec\-tures SIMD
  pour une architec\-ture cellulaire presque pure}

Nous avons d\'ej\`a \'evoqu\'e l'\'equivalence entre une fonction de transition
et une machine SIMD\@. %
La cellule d'espace est la m\'emoire locale \`a un PE, la table de
transition le PE, le r\'eseau d'interconnexion les liens entre les PE,
et la fonction d'application l'\'emulateur~(la machine) qui
d\'eplace les bits de la m\'emoire aux PE\@. %
Le choix d'une fonction d'application d\'etermine la taille de la table
de transition, et donc le nombre de machines virtuelles possibles
associ\'ees \`a cette fonction. %
Les contraintes qui p\`esent sur la conception d'une architecture
mat\'erielle r\'eelle ne sont pas les m\^emes que celles qui p\`esent
sur une architecture mat\'erielle virtuelle. %
Les machines virtuelles peuvent \^etre tr\`es sp\'ecifiques,
et pas n\'ecessairement universelles. %
Les machines SIMD r\'eelles doivent \^etre universelles. %
N\'eanmoins, le but de l'objet loi de transition, est bien de permettre
la sp\'ecification d'une architecture SIMD quelconque, et sa mise
en~{\oe}uvre la plus efficace possible. %

Que doit \^etre une fonction de transition pour simuler un plan de $N$~GAPP\@?

La cellule est compos\'ee %
des 3~registres~(C, NS et EW), %
des 3~bits de sortie de l'alu~(SM, CY et BW), %
des 2~bits de voisinage~(N ou S, et E ou W), %
et de 1~bit de ram~(M). %
La fonction d'application re\c{c}oit en plus des arguments usuels, %
les 11~bits de contr\^ole, %
les 7~bits d'adresse %
et un espace suppl\'ementaire de $N \times  128$~bits. %
Voir la figure~\mygetfig{gap}. %

Le travail de d\'ecodage des bits de contr\^ole est minimal et peut donc
\^etre sorti de la boucle d'application. %

Le bit de ram de la cellule est charg\'e depuis la m\'emoire.\\
Condition $\text{Ctl}_{8} \text{Ctl}_{9} \equiv 0$. %

\begin{verbatim}
*cellule &= CLEAR_RAM;
*cellule |= !!(ram[adresse_hi] & adresse_pos) << RAM;
\end{verbatim}

Le registre NS est charg\'e depuis la cellule au nord.\\
Condition $\text{Ctl}_{4} \text{Ctl}_{3} \text{Ctl}_{2} \equiv 2$. %

\begin{verbatim}
*cellule &= CLEAR_NS;
*cellule |= *north & NS;
\end{verbatim}

L'action existe aussi pour les cas $\text{NS} \leftarrow \text{S}$,
$\text{EW} \leftarrow \text{E}$ et $\text{EW} \leftarrow \text{W}$. %

L'index est compos\'e avec les 11~bits de contr\^ole et
les 7~bits d'\'etats de la cellule\footnote{%
  La table de transition a donc une taille de 1024~kilobytes,
  chaque byte contenant 7 bits utiles.}. %
Le nouvel \'etat de la cellule est calcul\'e. %

\begin{verbatim}
*cellule = transition[*cellule | controle];
\end{verbatim}

Le bit de ram est \'ecrit dans la m\'emoire.\\
Condition $\text{Ctl}_{8} \text{Ctl}_{9} \neq 0$. %

\begin{verbatim}
ram[adresse_hi] &= clear_ram;
ram[adresse_hi] |= !!(*cellule & GET_RAM) << adresse_pos;
\end{verbatim}

Les 3~conditions possibles pour l'acc\`es \`a la ram~%
(pas d'acc\`es ou read ou write) %
et les 5~conditions possibles pour l'acc\`es aux voisins~%
(pas d'acc\`es ou $\text{NS} \leftarrow (\text{N} \vee \text{S})$ %
ou $\text{EW} \leftarrow (\text{E} \vee \text{W})$) %
donnent 15 fonctions de boucle d'application, qui peuvent pr\'ecalculer
certaines valeurs associ\'ees aux conditions~(offset en bytes et en bits
de l'adresse par exemple), et dans le cas d'une instruction sans acc\`es
m\'emoire ni voisins, peut se ramener \`a %
\texttt{*cellule = transition[*cellule]}~%
(apr\`es un pr\'ecalcul \texttt{transition += contr\^ole}). %

On peut aussi ajouter le traitement du \emph{global output},
dont la condition de r\'ealisation est sous contr\^ole du s\'equenceur. %

\begin{verbatim}
global |= *cellule & NS;
\end{verbatim}

Nous avons tout de m\^eme besoin d'\'ecrire la fonction qui remplit la
table de transition, mais le passage par table simplifie
consid\'erablement la t\^ache, comme le montre clairement les
figures~\mygetcode{c-gapp-pe-h}, %
\mygetcode{c-gapp-pe-pe} et %
\mygetcode{c-gapp-pe-tbl}\footnote{%
  L'impl\'ementation de la construction de la table.
}. %

%%
%% OBJECTS/LOIS/COMPILE
%%
\subsection{Et leur compilateur:~la d\'ecompilation des lois pour des
  architectures SIMD}

Nous avons montr\'e comment un AC r\'ealis\'e par table
peut simuler une architecture SIMD\@. %
Comment transformer une fonction de transition en une s\'equence
d'instructions pour une machine SIMD\@? %
Soit une table de transition de $N$ bits de voisinage
et $R$ bits de r\'esultats. %
Construisons $R$ tables de taille $2^{N}$ bits. %
Nous cherchons un g\'en\'erateur d'instructions pour une machine SIMD pour
chacune des $R$ fonctions bool\'eennes associ\'ees \`a la fonction de
transition. %
Soit $T$ une table de $N$ bits de voisinage,
avec $R = 1$. %
La table $T$ d\'efinit un langage sur les $2^{N}$ index de la table. %
Le nombre de mots du langage \'etant fini, il est toujours possible de
construire un DFA pour reconna\^itre ce langage. %
Les expressions r\'eguli\`eres constituent donc un mode d'expression
possible d'une table de transition. %
Par exemple pour \arule{Life}, si le premier bit de l'index est le bit cible:\\
%%\texttt{(00*10*10*10*)|(10*10*10*10*)|(10*10*10*)}
$(00^*10^*10^*10^*)|(10^*10^*10^*10^*)|(10^*10^*10^*)$

\mytable{gapp-dfa}%
{Compilation d'un table pour le GAPP,}%
{Calcul d'une g\'en\'eration de \arule{Life}}%
{Une table repr\'esentant l'encodage du DFA de \arule{Life}, utilisant
  les noms des \'etats de la figure~\mygetfig{life-dfa}}

\mytable{gapp-dfa-optim}%
{et une optimisation imm\'ediate.}%
{Une optimisation de la table~\mygettbl{gapp-dfa}}%
{Il est possible d'obtenir un encodage plus \'economique\footnote{%
    En nombre d'instructions.} %
  du DFA de \arule{Life}, plus proche d'un encodage utilisant la
  nature arithm\'etique de la r\`egle. %
  L'optimisation pr\'esent\'ee repose quant \`a elle sur un m\'ecanisme
  g\'en\'eral.}

Chaque PE de la machine SIMD est d\'edi\'e \`a une cellule. %
La m\'emoire locale de chaque PE maintient l'\'etat de la cellule,
plus l'\'etat courant du DFA\@. %
Pour chacun des bits de voisinage de l'automate, chaque PE re\c{c}oit le
sous-ensemble du DFA construit \`a partir de l'ensemble des \'etats du DFA
pr\'ec\'edemment acc\'ed\'es. %
Le DFA est transmis sous la forme de triplets~(\'etat courant, \'etat si~1,
\'etat si~0). %
Pour finir, chaque PE re\c{c}oit la liste des \'etats acceptants. %

\clearpage

%%
%% OBJECTS/SEQUENCER
%%
\section{Les s\'equenceurs}

Nous pensions initialement que la phase de sp\'ecification de
l'architecture d'une exp\'erience cellulaire pouvait \^etre s\'epar\'ee de la
phase de s\'equencement d'ex\'ecution. %
Nous n'avons pas de solution g\'en\'erale satisfaisante au probl\`eme de
la s\'eparation de ces op\'erations, de la sp\'ecification\footnote{%
  Une fonction de s\'equencement exemplaire est la r\'egulation d'un temps
  relatif pour une trajectoire. %
  Certain AC pr\'esentent une \'evolution croissante tendant vers un point
  fixe ou un cycle limite court, passant par une phase de croissance
  tr\`es rapide suivie d'une phase de croissance tr\`es lente d'une
  variable macroscopique \'el\'ementaire comme la population d'un  plan de
  bits. %
  Le couplage de la variation de cette variable \`a la commande
  d'enregistrement du point de temps courant permet de r\'ealiser des
  films \`a \'echantillonnage variable et \`a variation de densit\'e
  spatiale constante.} %
et de l'impl\'ementation des s\'equenceurs. %
Les solutions particuli\`eres que nous avons envisag\'ees, test\'ees ou
impl\'ement\'ees reposent toutes sur une approche \aquote{classique}
centr\'ee sur des m\'ethodologies ou des outils d'impl\'ementation
d'interfaces. %
Cependant, nos exp\'eriences de traduction pratiques d'architectures
SIMD vers cellulaire, et cellulaire vers SIMD, nous am\`enent \`a proposer
une nouvelle approche des fonctions du s\'equenceur, bas\'ee sur la
d\'efinition d'une architecture mixte:~une architecture SIMD incluant des
primitives de traduction et de manipulation d'architectures
cellulaires pures. %
Le probl\`eme du s\'equenceur cellulaire devient alors celui, %
plus g\'en\'eral, de la sp\'ecification et du pilotage d'une machine SIMD
virtuelle.

\mydumpfull[t]{cantata-1}%
{}%
{Une session avec \asoft{Cantata}}%
{\asoft{Cantata} est le shell visuel interactif de \asoft{khoros}. %
  Dans cette session, nous utilisons les panneaux de contr\^ole de nos
  modules de g\'en\'eration de tables de transition et de g\'en\'eration de
  trajectoires cellulaires pour explorer l'espace des param\`etres
  statiques d'une r\`egle, via une visualisation tridimensionnelle de
  la trajectoire utilisant des outils de la biblioth\`eque
  \asoft{khoros}.}

\mytwodumpfull[t]{cantata-mk-rule}{cantata-cell}%
{}%
{Les panneaux de contr\^ole de nos modules \asoft{Cantata}}%
{La solution de \asoft{Cantata} aux probl\`emes de s\'equencement passe
  par l'existence de variables partag\'ees par les panneaux de contr\^oles
  des diff\'erents modules. %
  Un module de contr\^ole de boucle peut ainsi calculer un indice en
  fonction d'arguments de modules de g\'en\'erations d'espaces situ\'es en
  amont dans le graphe de production.}

\mycodetwocol[p]{tcl-cas}%
{}%
{D\'efinitions de proc\'edures de s\'equencements}%
{Ces proc\'edures travaillent sur des objets top-level destin\'es \`a
  l'interface humaine\footnote{%
    Soit par un tty, soit par une interface \asoft{Tk}. %
    \mygetdump[phrase]{casNew}}}

Nous avons cherch\'e les solutions possibles, exp\'eriment\'e plusieurs modes
de s\'equencement et impl\'ement\'e une interface \aquote{langage de commande} bas\'ee
sur~\asoft{Tcl} et une extension graphique interactive bas\'ee sur
\asoft{Tk}. %

La solution la plus simple consiste \`a b\^atir une commande sp\'ecialis\'ee
par classe de fonction de la librairie, puis \`a utiliser le shell Unix
comme s\'equenceur et le syst\`eme de fichiers comme syst\`eme de nommage des
objets interm\'ediaires. %
Cette solution, la solution Unix classique, reste irrempla\c{c}able. %
Nous l'avons largement utilis\'ee. %
Le sc\'enario de base est le suivant:
\begin{verbatim}
mk-rule-foo > foo.rule; mk-space-bar > bar.space
run-ca foo.rule bar.space | xdumper
run-ca -n 8192 foo.rule bar.space | do-stats > data
plot data
\end{verbatim}

Une autre solution classique consiste \`a d\'efinir un
langage~\cite[]{Almasi87,Eckart92b} dot\'e de structures de donn\'ees et de
contr\^oles adapt\'ees aux objets manipul\'es. %
Les outils de r\'ealisation existent\footnote{%
  lex, yacc, TXL~\cite[]{Cordy91}}, %
mais les choix des objets et des op\'erations primitives \'etant
pr\'ecis\'ement un des objets de notre recherche, ne peuvent venir avant
l'exp\'erimentation.

Une solution interm\'ediaire se trouve dans l'usage d'un interpr\`ete d'un
langage extensible permettant une manipulation interactive des objets
de la librairie. %
C'est la solution que nous avons retenue pour le workbench\footnote{%
  C'est aussi la solution retenue pour le contr\^ole de la machine
  CAM6\cite[]{Toffoli87}, sur base Forth.}, %
sur base \asoft{Tcl}. %

Nous avons \'egalement imagin\'e l'usage d'interpr\`etes sp\'ecialis\'es dans les
traitements num\'eriques, en particulier matriciels, comme support
d'interface \`a la librairie. %
En effet, nous avons r\'ealis\'e des manipulations de donn\'ees \`a l'aide
d'outils \`a la \asoft{matlab}:~\mysoftcite{Octave},
\mysoftcite{Rlab}. %
Ceux-ci constituent un envi\-ron\-nement bien adapt\'e \`a la manipulation
statistique des donn\'ees associ\'ees \`a une trajectoire, et sont par ailleurs
des langages suffisamment g\'en\'eraux pour la r\'ealisation de fonctions de
s\'equencement simples. %
L'id\'ee est d'ajouter les types de donn\'ees \aquote{espace} et
\aquote{r\`egle}, et une op\'eration d'application, ainsi que des
op\'erations de conversion \bialt{matrices}{espaces} et
\bialt{matrices}{r\`egles} \`a ces environnements.

La derni\`ere solution classique, est l'usage des outils de
visualisation scien\-tifique, en particulier de leurs shells graphiques
interactifs permettant la construction de graphes de traitement.
Nous avons utilis\'e \mysoftcite{Khoros}, dont l'outil de programmation
visuelle est \mysoftcite{cantata}. %
A fin de prototypage nous avons d\'evelopp\'e des panneaux de contr\^oles
\asoft{Khoros}, associ\'es \`a l'interface \aquote{commande Unix} de notre
librairie\footnote{%
  La solution la plus compl\`ete consisterait \`a d\'evelopper un toolkit
  \asoft{Khoros} bas\'e sur notre librairie.}. %
Le sc\'enario de base est alors une session interactive
\asoft{Cantata}, au cours de laquelle les sorties des modules
g\'en\'eration d'espaces de configurations et de g\'en\'eration de tables de
transitions, ad\'equatement param\'etr\'ees via les panneaux de contr\^ole
associ\'es, sont connect\'es aux entr\'ees correspondantes d'un module
\aquote{moteur cellulaire}, lequel est lui-m\^eme connect\'e \`a un module
de visualisation. %
L'int\'er\^et principal de cette approche, outre l'aspect visuel, r\'eside
dans la possibilit\'e, une fois l'adaptation aux normes d'interface
r\'ealis\'ee, d'utiliser la vaste biblioth\`eque de modules disponibles dans
ce genre d'environnement, et en particulier sous \asoft{Khoros}. %
Des modules de synth\`ese et manipulation ou d'analyse d'images de \asoft{Khoros}
peuvent \^etre utilis\'es en compl\'ement de nos modules. %

Aucune des solutions pr\'ec\'edentes ne nous a permis de nous affranchir
des probl\`emes de s\'equencement. %

La solution que nous proposons pour am\'eliorer cette situation passe
par la mise en pratique d'une association architecturale mixte
\bialt{SIMD}{cellulaire}. %

Nous pensons que l'association d'AC localement universels~(PE) \`a des
AC non localement universels fournit un cadre conceptuel adapt\'e au
calcul cellulaire appliqu\'e \`a l'\'etude des AC dans un contexte mat\'eriel
s\'equentiel ou de parall\'elisme non massif. %

En compl\'ement de la transformation d'un AC donn\'e par sa table de
transition en s\'equence d'instructions pour une machine SIMD, nous
consid\'erons la transformation d'un algorithme cellulaire en un
ensemble d'AC donn\'e par leur table, associ\'e \`a une s\'equence
d'instructions pour une machine SIMD, r\'ealisant le s\'equencement des
sous-AC de l'algorithme.

Le probl\`eme du s\'equencement devient alors un programme d'\'etude des
architectures de calcul cellulaire appliqu\'e \`a l'exp\'erimentation et
l'\'etude des AC\footnote{%
  Les architectures cellulaires sont \'evidemment applicables \`a d'autres
  domaines que les AC:~%
  traitement d'image, visualisation, certain calculs num\'eriques,
  simulation d'architecture mat\'erielle, en g\'en\'eral tous les domaines
  pour lequel le SIMD est applicable.}. %
Nous insistons sur le point suivant:~%
une architecture cellulaire n'a pas besoin d'une machine parall\`ele.
Hypoth\`ese:~\aquote{On peut impl\'ementer de mani\`ere effi\-cace une
  architecture cellulaire sur une machine s\'equentielle}. %
Une telle impl\'ementation a pour vocation de supporter des
impl\'ementations cellulaires d'algorithmes cellulaires sur machine non
cellulaire sans perte d'efficacit\'e\footnote{%
  Les algorithmes du traitement d'image sont candidats \`a la validation
  de cette hypoth\`ese.}. %
Les avantages sont multiples:~coh\'erence et clart\'e de la description
des algorithmes dans leur environnement architectural naturel. %
Simplicit\'e\footnote{%
  Au sens o\`u la nature cellulaire de l'algorithme n'a pas \`a \^etre
  traduite en boucles.}, %
facilit\'e de debug,
portabilit\'e et efficacit\'e de la programmation par look-up table, %
enfin parall\'elisation facile, ou plut\^ot, mapping sur une architecture
parall\`ele arbitraire d'une description cellulaire d'un algorithme.

Les AC peuvent alors \^etre vus comme les primitives d'une
architecture cellulaire, et ceux-ci deviennent ainsi les moyens
d'\'etude de leur propre univers, car ils posent de probl\`emes de
visualisation, de r\'eduction num\'erique, de sp\'ecification d'architecture
mat\'erielle et de dynamique cellulaire, toutes choses pour lesquelles
le calcul cellulaire est applicable.

%%
%% Local Variables:
%% mode: latex
%% mode: auto-fill
%% TeX-master: "top"
%% TeX-command-default: "mytex"
%% iso-accents-minor-mode: nil
%% Local IspellParsing: "latex-mode ~latin1"
%% Local IspellPersDict: "~/.ispell_lisp"
%% LocalWords: "d'AC l'AC s\'equenceur SIMD binding concat\'enation bytes byte thy"
%% LocalWords: "GET-REF nombre-de-voisin m\'egabytes combine-func l'isotropie SETR"
%% LocalWords: "mapping p\'enalisante cross-over multiplexage BIG loop Langton"
%% LocalWords: "func apply apply fig-apply pipeliner Moore shift quote cellsim"
%% LocalWords: "ref-mem-struct adressables br bp von Neuman \'economiseurs pattern"
%% LocalWords: "rk log times Ctl sub Ctl sub NS EW EW auto-reproducteur gapp"
%% LocalWords: "paren lambda auto-re uvre mem-struct indirection timer sum XOR"
%% LocalWords: "cycliquement time-xor s\'equencement seq-vs-pipe accumulatrice ACU"
%% LocalWords: "ACU SISD static-param timers l'applicatif trap PE l'\'emulateur SM"
%% LocalWords: "gap-pe gap l'alu CY BW ram CLEAR hi pos north kilobyte controle"
%% LocalWords: "clear GET read write pr\'ecalculer pr\'ecalcul output DFA life-dfa"
%% LocalWords: "life-tree state Center South East West NW SW lookup-table fonc"
%% LocalWords: "lisation tion Xor s\'e paration Neumann architec ture kilobytes Tk"
%% LocalWords: "d\'ecompilation triplets acceptants gapp-dfa gapp-dfa-optim shell"
%% LocalWords: "s\'equenceurs transmutations transmutation cantata khoros Tcl lex"
%% LocalWords: "cantata-mk-rule cantata-cell nommage mk-rule-foo foo.rule run-ca"
%% LocalWords: "mk-space-bar bar.space xdumper do-stats yacc TXL workbench CAM"
%% LocalWords: "Forth d'interpr\`ete MathLab Rlab shells prototypage toolkit debug"
%% LocalWords: "sous-AC impl\'ementation impl\'ementations portabilit\'e look-up tty"
%% LocalWords: "tcl-cas s\'equencements top-level casNew matriciel c-gapp-pe-h"
%% LocalWords: "Field c-gapp-pe-pe purements c-gapp-pe-tbl implode explode in"
%% LocalWords: "bit-packed bit-expanded inline c-gapp-arch-h extenso c-gapp-arch"
%% LocalWords: "c-gapp-arch-alu l'ALU"
%% End:
